\begin{frame}{Fibonacci Convention}
\[F(0)=0,\qquad F(1)=1,\qquad F(n)=F(n-1)+F(n-2)\quad(n\ge2)\]
\begin{block}{한 강의, 한 convention}수식, pseudocode, C/Java 코드 모두 0-based state $F(0)$부터 사용한다.\end{block}
\end{frame}
\begin{frame}{Fibonacci를 DP Template로 표현하면}
\small
\dpdesign
{$F(i)$: $i$번째 Fibonacci 수}
{$F(i)=F(i-1)+F(i-2)$}
{$F(0)=0,\ F(1)=1$}
{$i=2,3,\ldots,n$ 또는 top-down dependency 호출}
{$F(n)$}
\end{frame}
\begin{frame}{Naive Recursive Fibonacci}
\small\begin{algorithmic}[1]
\Procedure{FibRecursive}{$n$}
\If{$n\le1$}\State\Return $n$\EndIf
\State\Return \Call{FibRecursive}{$n-1$}+\Call{FibRecursive}{$n-2$}
\EndProcedure
\end{algorithmic}
\begin{itemize}\item base case는 종료와 정의를 동시에 담당\item 두 branch가 독립처럼 보이지만 같은 state를 다시 요청\end{itemize}
\end{frame}
\begin{frame}{\texorpdfstring{$F(6)$}{F(6)} Partial Call Tree}
\centering\begin{tikzpicture}[level distance=10mm,
level 1/.style={sibling distance=58mm},level 2/.style={sibling distance=28mm},
level 3/.style={sibling distance=14mm},every node/.style={tree node,minimum width=13mm,font=\small}]
\node{$F(6)$}
 child{node{$F(5)$} child{node[dp repeat]{$F(4)$} child{node[dp repeat]{$F(3)$}} child{node[dp repeat]{$F(2)$}}} child{node[dp repeat]{$F(3)$} child{node[dp repeat]{$F(2)$}} child{node{$F(1)$}}}}
 child{node[dp repeat]{$F(4)$} child{node[dp repeat]{$F(3)$} child{node{$\cdots$}} child{node{$\cdots$}}} child{node[dp repeat]{$F(2)$}}};
\node[callout]at(0,-4.2){partial call tree: 같은 입력이지만 서로 다른 호출과 activation record};
\end{tikzpicture}
\end{frame}
\begin{frame}{중복 호출을 세어 보자}
\begin{columns}[T]\column{.48\textwidth}
\begin{block}{반복 state}$F(4)$, $F(3)$, $F(2)$가 서로 다른 tree node로 여러 번 등장한다.\end{block}
\column{.48\textwidth}
\begin{block}{Tree와 DAG}call tree의 node는 호출 사건이다. 동일 state를 합치면 dependency DAG가 된다.\end{block}
\end{columns}
\dptakeaway{DP는 tree의 같은 label을 하나의 state로 합쳐 한 번만 계산한다.}
\end{frame}
\begin{frame}{Naive Recursion의 비용}
\[
T(n)=T(n-1)+T(n-2)+\Theta(1)
\]
\begin{itemize}\item 시간: $O(2^n)$의 단순 upper bound; 더 정밀하게 $\Theta(\varphi^n)$\item recursion depth: $\Theta(n)$\item state 종류는 $n+1$개뿐인데 호출 node는 exponential\end{itemize}
\end{frame}
\begin{frame}{Fibonacci Checkpoint}
\begin{enumerate}\item $F(5)$와 $F(3)$ branch가 공유하는 state는?\item 저장해야 할 서로 다른 state 수는?\item recursion depth는?\end{enumerate}
\begin{exampleblock}{답}표시된 partial tree에서는 대표적으로 $F(3)$과 그 descendant states가 중복된다. 전체 tree에서는 $F(3),F(2),F(1),F(0)$이 다시 요청될 수 있다. 서로 다른 state는 $F(0)..F(n)$의 $n+1$개이고, depth는 $\Theta(n)$이다.\end{exampleblock}
\vfill\muted{같은 입력의 답을 저장하면 이후 호출은 그 subtree를 다시 확장하지 않는다.}
\end{frame}
